
-- 必要 -----------------------------------------------------------




----------------------------------------------------------------

しかし現状として，カラオケを通じた自己開示に関する先行研究は存在せず，これらの要素が人間関係の親密化にどのように寄与するかは十分に明らかにされていない．特に，歌唱を伴う自己開示はより内面的な要素を含むため，その基礎となる音楽嗜好の共有を通じた自己開示と受容の関係を明らかにする必要がある．
まず音楽嗜好を通じた自己開示そのものが受容性に影響する要因がまだ明らかではない．本研究では，カラオケにおいて段階的に自己開示を促進することにより，受容性が向上するかどうかの検証を行う．

\section{音楽の好みの類似性（未完成）}
カラオケ場面における音楽の好みの類似性は，受容に大きく影響すると考えられる．本節ではまず，自己開示や受容にはどのような類似性が影響するかを述べる．そして，音楽の好みの類似性がどのように構成され，それが自己開示に与える影響について述べる．

\subsection{類似性と受容の関係}
類似性ー対人魅力の話とか

類似性により，共感反応を示しやすい → 受容が促されるって話


\subsection{音楽の好みの類似性と受容の関係}
音楽の好みの類似性の構成要素（ジャンル，アーティスト，曲）

音楽の好みの類似性は，複数の要素から構成される．
アーティストの類似性は，同一アーティストの曲を好むか，あるいは似たような特徴を持つアーティストを好むか
ジャンルの類似性は，J-POP，ロック，アニメソングなど，大まかな音楽カテゴリーの一致度
楽曲レベルでは，テンポやキー，メロディーライン，歌詞の内容など，より詳細な音楽的特徴の類似性

例えば，好みのアーティストや曲が一致している場合，その共通点を起点として会話が発展しやすい．また，相手と音楽の好みが似ていることを認識することで，開示ハードルが下がる．相手への魅力が向上

システム実装の観点からは，アーティスト名やジャンル，テンポやメロディーラインといった要素は比較的扱いやすい．
一方で，楽曲の詳細な特徴や，アーティストや曲に対する主観的な好みの度合いは定量化が困難であり，システムによる支援には工夫が必要となる．


以上より，カラオケにおける音楽の好みの類似性は，xxな理由から，受容を促すと考えられる．

次節では，これらを踏まえ，カラオケ場面における自己開示支援システムの要件について述べる．

\begin{comment}

\section{自己開示における趣味の位置づけ}
% 趣味の位置づけ・特徴
自己開示の研究において，趣味に関する開示は一般的に表層的なものとして位置づけられてきた．1章で取り上げた丹羽・丸野\cite{自己開示尺度}の自己開示尺度では，趣味・趣向に関する開示は最も浅いレベルⅠに位置づけられ，「一般的な自分の好みについての情報で，特に社会的常識から逸脱する趣味ではない限り，人格を疑われたりするものではない」と説明されている．
しかし，趣味の開示には段階的な自己開示を実現する上で重要な可能性を持つ以下の特徴がある．


第1に，開示のしやすさは，初期段階における自己開示の促進に寄与する．相手の反応を予測しやすい話題から始めることで，より深い自己開示への心理的ハードルを下げることができる．
第2に，個人的意味の多層性は，社会的浸透理論が示す段階的な関係性の発展と整合する．表層的な好みの共有から始まり，徐々により個人的な意味や価値観の共有へと発展させることが可能である．
第3に，類似性の認識のしやすさは，受容性の向上に寄与する．共通点を見出しやすい話題から始めることで，相互理解と信頼関係の基盤を構築することができる．
% エビデンスなし
これらの特徴は，本研究が目指す段階的な自己開示の実現において重要な意味を持つ．

\end{comment}

\begin{comment}
% ＜カラオケの理由＞
通常の会話での音楽の好みの開示と異なり，カラオケでは実際の音楽が流れ，その場で歌を披露するという直接的な体験を伴う．この特徴により，カラオケは以下のような独自の利点を持つ自己開示の場として機能する．


単に会話内で音楽の好みを開示するのに対し，実際に音楽が流れ，相手の前で直接歌を披露するとなったとき，「上手く歌えなかったらどうしようか」「これは相手に受け入れられるような曲だろうか」といった，失敗不安からくる否定的感情が湧くものである．このような心理的なリスクを伴う自己開示は，より深い内面性を持つ開示として捉えることができる．

% ＜自己開示における心理的リスクと内面性の関係性＞
さらに，自己開示の相互性を研究する場としても，カラオケは優れた特徴を有している．通常の会話場面では相手の話を遮る，話題を変えるといった現象が発生し，一方的に自己開示をし続けたり，自己開示の割合に大きな偏りが生じたりする可能性がある．そのため，（AとBという2者関係において，）Aが自己開示をしてBが同等の自己開示を返すという相互作用を明示的に行うことは困難である．

% ＜通常の会話場面における自己開示の課題＞
一方，カラオケでは「歌を歌う」「相手の歌を聴く」という役割が明確に分かれており，さらにそれを交互に繰り返す状況を容易に作り出すことができる．これにより，自然な相互作用の空間を再現しつつ，自己開示の相互性を明示的に観察することが可能となる．
また，カラオケは単に1人で歌を楽しむものではなく，自分が歌い，相手が聴くという空間であることを意識せざるを得ない．数家\cite{数家鉄治_2008}は，カラオケは共同行為に必要な自己と他者の相互依存性をより緊密にできる場と述べている．



% 言語的なものじゃないけど，まー良いでしょの精神．
カラオケにおける自己開示は言語的なものではないが，選曲・歌唱は，個人の内面や特性を他者に開示するという本質的な機能を持っていると考えられる．また，選曲・歌唱を通じた会話となれば言語的なものであり，本来の自己開示と同様に扱うことが可能であると考えた．このように，音楽嗜好を通じた自己開示，特にカラオケにおける自己開示は，段階的な親密化のプロセスを理解し支援するための重要な研究対象となる．

音楽の好みを通じた自己開示では，その深さに応じて異なる段階が存在すると考えられる．このような段階性を明確にするため，次項では社会的浸透理論の枠組みに基づき，カラオケにおける自己開示の段階を体系的に整理する．


% 表層の音楽の好みに着目した利点
音楽の好みという比較的開示しやすい情報から始めることで，心理的安全性を確保しながら段階的な自己開示を促進することができる．
また，音楽の好みは，アーティスト，ジャンル，楽曲などの具体的な要素を通じて類似性を判断しやすい．これにより，システムによる支援が特に効果的に機能すると考えられる．さらに，表層的な好みの共有では，相手の反応や受容の度合いを比較的客観的に観察することが可能である．これは支援システムの効果を評価する上で重要な利点となる．
\end{comment}



カラオケの自己開示4段階
\begin{enumerate}
    \item 表層：音楽の好み，好き嫌い
    \item 中間層：場の空気を読んだ選曲判断，他者との関係性を考慮した曲選び
    \item 内層：音楽に対する個人的な価値観，特定の曲に対する思い入れ，失敗への不安/成功への期待
    \item 核：極めて個人的な経験との結びつき
\end{enumerate}
